% ==========================================================================
% Chapter 4: Social and Environmental influences
%
% %: percentage required by the EEE 4000 manual (General Instruction 7d). Graded under
% CO6/PO(f) (societal, health, safety, legal and cultural issues) and
% CO7/PO(g) (environmental impacts/sustainability).
% ==========================================================================

\chapter{Social and Environmental Influence}
\label{ch:social}

This chapter discusses the broader social and environmental implications of
the automated ECG arrhythmia classifier developed in this thesis. The work
sits at the intersection of biomedical engineering, deep learning, and
resource-constrained embedded systems, and its impact is therefore assessed
along three axes: the people who use it or are monitored by it
(Section~\ref{sec:societal}), the environment in which it is manufactured
and operated (Section~\ref{sec:environment}), and the long-term
sustainability of the solution (Section~\ref{sec:sustainability}).

\section{Societal Impact}
\label{sec:societal}

The direct beneficiaries of this work are cardiac patients and the
clinicians and caregivers who monitor them. The principal societal
contribution is an automated, microcontroller-based approach to
democratising arrhythmia screening: a small, wearable device that can
deliver continuous monitoring at low cost to people who currently cannot
afford it. In Bangladesh, cardiologists are concentrated in urban centres,
so the patient-to-doctor ratio outside major cities is unfavourable and
tertiary care is often unavailable. A device of this kind could change how
dangerous arrhythmias, such as ventricular ectopy, are detected, moving
detection from the hospital to the home, and catching conditions on the
patient's own time rather than only during scheduled visits.

\subsection{Financing and Health Effects}
\label{subsec:financial_health}

\textbf{Financial:} The proposed system is deliberately designed around
low-cost components. The microcontroller used in this work, the STM32,
costs only a few dollars, and the deployed model has a total footprint of
approximately 288~KB of flash with no dependence on cloud infrastructure.
The processing and inference cost is therefore a one-time hardware
expense rather than a recurring data or servicing fee. By contrast, the
conventional strategy of Holter monitoring followed by manual cardiologist
review has been shown to cost considerably more per patient and does not
scale to continuous monitoring of the general population. In the
Bangladeshi context, where the cost of cardiac care is often a barrier, an
affordable wearable monitor is one way to help address this problem. A
classifier of this kind could ease the financial burden that arrhythmia
places on patients, reducing the morbidity and mortality costs, as well
as the costly hospitalisation for stroke and heart failure, that
undetected arrhythmias can cause.

\textbf{Health:} As a biomedical work, the primary impact on health is
tied to classification performance, specifically, the class of missed
beats and the clinical risk each represents. With a testing sensitivity
of 92.79\% for the highest-risk class, and, as Section~\ref{subsec:ec57}
shows, a false-positive rate of only 0.73\% of samples, the system results
in fewer missed life-threatening beats. Early detection of atrial and
other arrhythmic events, such as fibrillation-like episodes, can enable
earlier treatment and help avoid strokes and sudden cardiac death. The
device is non-invasive: it senses the ECG through skin electrodes only, so
it exposes the patient to no ionising radiation and no physical risk
beyond that of normal electrode placement. The main health considerations
for extended wear are skin irritation from the electrodes and, as with any
battery-operated wearable, thermal and chemical battery safety, both of
which are addressed through the use of certified, low-power components.

\subsection{Safety and Legal/Cultural Issues}
\label{subsec:safety_legal_cultural}

\textbf{Safety:} As a low-voltage, battery-powered system, the electrical
hazards involved are small compared with mains-powered medical electronic
equipment. The main hazards are data loss from a missed classification on
an unattended monitor, and alarm fatigue, in which too many false alarms
degrade the caregiver's ability to respond appropriately. Both are
addressed in the design: the classifier is evaluated statistically using
the false-positive rate per class (Section~\ref{subsec:ec57}), and the
recall-floor analysis of Section~\ref{subsec:roc} quantifies the precision
cost of any mandated detection rate, so that a clinician can select an
operating point that bounds the false-alarm burden before the device is
deployed.

Were this system used as a medical decision-support device, it would need
to comply with the relevant regulations governing complex computer
programs used in patient care. The classification methodology is already
designed to be compliant with the testing and reporting standards for
cardiac rhythm algorithms~\cite{aami_ec57_2012}. Deployed as a regulated
medical device in Bangladesh, it would need to obtain clearance from the
Directorate General of Drug Administration (DGDA), whose guidelines are
broadly similar to international frameworks such as the FDA and CE
marking routes. Data protection is an additional consideration: although
the model in this work was trained on de-identified public data from
PhysioNet, a deployed wearable would handle identifiable patient
information and must comply with the applicable data-protection regime,
including informed consent and secure storage.

\textbf{Cultural and ethical:} The most significant ethical
consideration concerns the attribution maps generated by the model
(Subsection~\ref{subsec:explain}). While plausible, these maps were not
found to be reliably faithful under deletion-curve testing and should
therefore be presented to clinicians as qualitative, supporting evidence
rather than as causal explanations. An ethically sound deployment of the
model should accordingly position it as a screening aid that keeps the
cardiologist in the loop, rather than as an autonomous diagnostic tool.
On equity, the device is low-cost, fully battery-powered, and operates
entirely on-device, so it can be used without internet connectivity in
rural or low-income settings, offering better equity of access than
cloud-based alternatives. However, over-reliance on automated screening
carries a real risk of reducing the number of trained ECG interpreters
over the long run; the effects on the employment of ECG technicians
should be monitored as such devices become more widespread.

\section{Impact on Environment}
\label{sec:environment}

The full life cycle of the device should be considered when assessing its
environmental impact.

\textbf{Manufacturing:} The hardware (a board containing electrodes and
a generic microcontroller) has a material footprint dominated by the
printed circuit board, the battery, and the enclosure, all components
whose manufacturing impact is well understood. The largest computational
cost is incurred during model development: training the deep learning
models required a number of GPU-hours. This cost is incurred only once,
at design time, and is offset by the fact that the deployed system
performs inference in short bursts only: the microcontroller draws power
for 185~ms per beat classification (Section~\ref{subsec:mcu}) and idles in
between, which is well suited to intermittent use in a battery-powered
wearable.

\textbf{Operation:} Compared with a policy of hospital-based monitoring,
a wearable classifier that substitutes for Holter monitoring and routine
hospital visits requires far less energy: low power draw, small size, and
on-device computation replace the travel, hospital energy use, and
clinician time associated with in-person monitoring. Every hospitalisation
avoided also represents a saving in transportation fuel. The operating-phase
footprint of the solution is consequently very low compared with the
status quo, and its net effect on emissions is likely negative.

\textbf{End of Life:} As an electronic device, the unit contributes to
Bangladesh's growing e-waste problem, though a small, microcontroller-based
wearable represents a smaller burden than larger hospital monitoring
equipment. Both the device and its battery need to be recycled and
disposed of through appropriate channels rather than household trash.
Because the classifier is implemented in software, it can be updated to
extend the device's useful life without requiring new hardware. The
quantised INT8 model is also very compact, at 29.9~KB, allowing it to run
for years on inexpensive, existing hardware, lowering the rate of
device replacement and the resulting e-waste.

\section{Sustainability Issues}
\label{sec:sustainability}

The sustainability of the solution rests on four factors. First,
\emph{spare parts and maintenance}: the design is based on commodity
microcontroller hardware that is readily available in Bangladesh, and
because all logic resides in firmware, a failed unit can be replaced
cheaply and patched locally rather than depending on a proprietary
cloud-based service that could be discontinued. Second, \emph{service
life}: the STM32 hardware platform has a long production lifecycle, and
the software footprint is small enough that the model can be updated
in place, allowing the device to remain in service for many years without
requiring hardware changes. Third, \emph{scalability}: both the per-device
cost and the per-inference energy consumption are low enough that the
solution could scale from a single clinic pilot to population-level
screening without a proportionate increase in running cost; the larger
challenge to scaling is the clinical workforce available to follow up on
screening results. Fourth, \emph{import dependence}: most of the core
hardware, including the highest-value components, is imported, but the
ability to train and deploy the model is a domestic capability, and it can
be reproduced, retrained on local data, and improved without foreign
supervision, which helps offset this dependence.

This work also advances several United Nations Sustainable Development
Goals. Its greatest contribution is to SDG~3 (Good Health and Well-being),
through the goal of reducing premature mortality from non-communicable
diseases such as cardiovascular disease by one third by 2030, which
continuous, affordable, out-of-hospital arrhythmia screening directly
supports. It also supports SDG~9 (Industry, Innovation and Infrastructure)
by building embedded machine learning research and engineering capability
in Bangladesh, and SDG~12 (Responsible Consumption and Production) by
offering a small, long-lived, low-energy device as an alternative to
disposable or energy-intensive equipment.
